\section{Arena Design}
\label{sec:arena}

The arena is a model of the task and does not reproduce the competition field. It keeps the structure of the task: a line to follow, then a
room in which coins are collected at one wall and delivered, by colour, to another. It was
designed together with the robot model, and most of its dimensions follow from two numbers
of \cref{sec:mechanical}: the pitch of the infrared array and the blind zone of the camera.

\subsection{Layout}
The field measures $3.6\times2.4\unit{m}$ and is enclosed by walls 0.3\unit{m} high
(\cref{fig:arena}). A partition with a doorway 0.8\unit{m} wide divides it into an entry
lane and a room. The robot starts in the lane on a black line, which leads through an S-bend
of two arcs of 0.6\unit{m} radius, through the doorway, to a black start square. Three pick
bays lie along the left wall and three drop zones along the front wall. Up to three boxes of
$100\times100\times150\unit{mm}$ stand in the room as obstacles.

\begin{figure}[htbp]
  \centering
  \begin{tikzpicture}[x=3.3cm, y=3.3cm]
    \arenaplan
    \arenacoins
    \arenaobstacles{1.90/0.55, 1.75/-0.15}
    % robot at the start, heading +x
    \draw[draw=ink, fill=panel] (-0.017,-0.075) rectangle (0.163,0.075);
    \node[note, anchor=north] at (0.07,-0.09) {start};
    % station poses
    \foreach \bx in {1.5,2.0,2.5} { \draw[->, cA, line width=0.8pt] (\bx,0.93) -- (\bx,1.07); }
    \foreach \dy in {0.7,0.2,-0.3} { \draw[->, cA, line width=0.8pt] (2.33,\dy) -- (2.47,\dy); }
    % labels
    \node[note] at (0.3,1.55) {entry lane};
    \node[note] at (1.2,-0.4) {room};
    \node[note, anchor=south] at (2.0,1.83) {pick bays 1 to 3, QR code on the wall};
    \node[note, rotate=-90, anchor=south] at (3.34,0.2) {drop zones 1 to 3, QR code on the wall};
    \node[note, anchor=north] at (1.46,0.52) {start square};
    \node[note, anchor=west] at (0.92,0.9) {doorway};
    \node[note, text=cB, anchor=west] at (1.97,0.52) {obstacle};
    \node[note, text=cA, anchor=east] at (2.3,0.83) {station pose};
    % axes and scale
    \draw[->] (-0.2,-0.5) -- (0.1,-0.5) node[note, right] {$x$};
    \draw[->] (-0.2,-0.5) -- (-0.2,-0.2) node[note, above] {$y$};
    \draw[|-|] (2.3,-0.72) -- (3.3,-0.72);
    \node[note, below] at (2.8,-0.73) {1\,m};
  \end{tikzpicture}
  \caption{Plan of the arena, to scale, with obstacle preset a. The arrows are the station
  poses at which the navigation hands over to the line follower.}
  \label{fig:arena}
\end{figure}

\subsection{Floor Markings Dimensioned from the Array}
\label{sec:markings}
The infrared array has eight channels at a pitch of $p=15\unit{mm}$ and a span of
105\unit{mm}. Three kinds of marking are derived from these two numbers (\cref{fig:markings}).
\begin{description}
  \item[Lines are 12\unit{mm} wide.] A line narrower than one pitch can never cover two
        channels completely. At most two neighbouring channels respond to it, each
        partially, and the ratio of their readings gives the position of the line between
        them.
  \item[Markers are 120\unit{mm} wide.] The start square and the crossbars of the pick
        bays are wider than the whole array, so that all eight channels go dark together
        when the array reaches one. Nothing else on the floor can produce that pattern, and
        a line, which darkens two channels at most, cannot be mistaken for it.
  \item[Drop squares are 250\unit{mm} wide.] The robot releases a coin when the array
        reaches the near edge of the square. The coin then lies 0.13\unit{m} further in,
        which is 5\unit{mm} from the centre of a square of this size.
\end{description}

\begin{figure}[htbp]
  \centering
  \begin{tikzpicture}[scale=0.58]
    % units: 10 mm. Left: array over a line. Right: array over a crossbar.
    \begin{scope}
      \fill[black] (0.9,-2.2) rectangle (2.1,2.2);
      \foreach \i in {0,...,7} {
        \draw[draw=ink, fill=white] ({(\i-3.5)*1.5},0) circle (0.4);
      }
      \draw[draw=cB, fill=cB!35] (0.75,0) circle (0.4);
      \draw[draw=cB, fill=cB!35] (2.25,0) circle (0.4);
      \draw[|-|] (-5.25,-1.3) -- (5.25,-1.3);
      \node[note, below] at (-2.5,-1.35) {span 105\,mm};
      \draw[|-|] (-5.25,1.0) -- (-3.75,1.0);
      \node[note, above] at (-4.5,1.05) {15\,mm};
      \node[note, above] at (1.5,2.25) {line, 12\,mm};
    \end{scope}
    \begin{scope}[xshift=14.5cm]
      \fill[black] (-6,-0.6) rectangle (6,0.6);
      \fill[black] (-0.6,-2.6) rectangle (0.6,-0.6);
      \foreach \i in {0,...,7} {
        \draw[draw=cB, fill=cB!55] ({(\i-3.5)*1.5},0) circle (0.4);
      }
      \draw[|-|] (-6,1.3) -- (6,1.3);
      \node[note, above] at (0,1.35) {crossbar, 120\,mm};
      \node[note, anchor=west] at (0.8,-1.8) {stem};
    \end{scope}
  \end{tikzpicture}
  \caption{The infrared array over a line and over a marker, to scale. The circles are the
  spots of the eight channels. A line darkens two channels at most, a marker all eight.}
  \label{fig:markings}
\end{figure}

\subsection{Stations Placed for the Blind Zone}
\label{sec:stations}
The camera does not see the last 0.15\unit{m} of floor (\cref{eq:blind}). A pick bay is
therefore a T on the floor: a stem 0.30\unit{m} long that leads to a crossbar, with the coin
0.13\unit{m} beyond the crossbar. The distance of 0.13\unit{m} is the distance from the
array to the space between the open jaws, so that a robot which stops with its array on the
crossbar has the coin between its jaws. A drop zone is a stem of the same length that ends
in the drop square.

The navigation brings the robot to a station pose on the extension of the stem and short of
its foot, 0.15\unit{m} for the bays and 0.175\unit{m} for the drop zones, and facing along
it. From there the line follower drives straight until the array finds the stem
(\cref{sec:docking}). The length of the stem is what absorbs the error of the navigation: a
lateral error of a few centimetres at the foot is steered out before the crossbar is
reached.

\subsection{QR Codes and Wall Textures}
A QR code of 160\unit{mm}, centred 0.11\unit{m} above the floor, is mounted on the wall
behind every station. It encodes the name of a colour. At a pick bay it names the colour of
the coin, and at a drop zone the colour that the zone accepts. The codes are generated with
the highest level of error correction and a wide quiet zone, so that they remain readable
when seen at an angle. The assignment of colours to bays and to zones is a launch argument,
which makes each of the 36 combinations of pick order and drop order a different mission.

Plain walls have no visual features, and a visual SLAM system cannot localise against them.
The walls therefore carry generated textures with dense, irregular detail. They exist only
for this purpose.

\subsection{One Source for the Dimensions}
The world is written as a parametrised file whose dimensions are defined once at its head.
Two other files repeat a few of them: the list of station poses read by the mission, and
the positions of the drop squares read by the referee. Both name the parameters that they
mirror. A change of the arena therefore requires edits in three places, which is a known
weakness of the present arrangement.
